\documentclass[twocolumn]{article}
\begin{document}
\section*{Starred forms}
\begin{figure*}[t]
  \centering
  \begin{subfigure}[b]{0.48\textwidth}
    \includegraphics[width=\linewidth]{wide/left}
    \caption{Left half.}\label{fig:wide-a}
  \end{subfigure}
  \begin{subfigure}[b]{0.48\textwidth}
    \includegraphics[width=\linewidth]{wide/right}
    \caption{Right half.}\label{fig:wide-b}
  \end{subfigure}
  \caption{A figure that spans both columns.}
  \label{fig:wide}
\end{figure*}

\begin{figure}
  \includegraphics{narrow}
  \caption{A figure in one column.}\label{fig:narrow}
\end{figure}

\begin{table*}
  \caption{A table that spans both columns.}\label{tab:wide}
  \begin{tabular}{lr}
    Name & Value \\
    alpha & 1 \\
  \end{tabular}
\end{table*}

\begin{equation}
  a = b \label{eq:plain}
\end{equation}
\begin{equation*}
  c = d \label{eq:starred}
\end{equation*}
\begin{align*}
  e &= f \\
  g &= h
\end{align*}
\begin{flalign}
  i &= j &
\end{flalign}

Figures~\ref{fig:wide}, \ref{fig:wide-b} and \ref{fig:narrow}; Table~\ref{tab:wide};
Equation~\eqref{eq:plain}.
\end{document}
